\begin{textsample}{POS Dim 2 – rn – Score 20.00 – rn/rn\_inf\_41.txt}  \label{ex:f2_pos_210}
Conforme a LDB ( BRASIL, 1996 ) em seu \textbf{artigo} 11 : “ Os Municípios incumbir -se-ão de : [... ] Autorizar, credenciar e supervisionar os estabelecimentos do seu sistema de ensino ”.
Segundo levantamento feito pela UNDIME, 132 \textbf{municípios} do Rio Grande do Norte possuem o Conselho Municipal de Educação, no entanto, apenas 36 possuem Sistema Municipal de Ensino.
Nesse caso, a maioria dos \textbf{municípios} \textbf{permanece} integrada ao Sistema Estadual de Ensino.
É importante a busca por garantir maior autonomia dos Conselhos \textbf{Municipais} de Educação, assim como, não se esgota a possibilidade de desenvolver \textbf{políticas} \textbf{municipais} de educação infantil e consequentemente, sua \textbf{normatização} local-municipal, mesmo que ainda não esteja consolidado um Sistema próprio de ensino.
Conforme aponta Nunes, Corsino e Didonet ( 2011 ), as normas referentes à \textbf{regulamentação} compreendem orientações e exigências que demarcam o funcionamento das \textbf{instituições} com vista a contribuir com a implementação da \textbf{política} \textbf{municipal} de educação infantil.
Orienta -se que os \textbf{municípios} tenham \textbf{regulamentações} acerca da organização e funcionamento de \textbf{instituições} públicas e privadas de educação infantil, acompanhadas pelo Conselho Municipal de Educação.
Tais normativas podem definir, entre outros aspectos : critérios para organização de grupos de crianças nas \textbf{instituições} e a relação entre quantidade de criança e de professores em cada grupo ; processos de admissão e de formação continuada de profissionais da educação infantil ; formas de \textbf{gestão} ; organização de espaços pedagógicos e definição de padrões mínimos de infraestrutura ; e orientações gerais relativas às suas propostas pedagógicas.
Essas definições precisam considerar as especificidades de cada \textbf{município} e/ou \textbf{instituição} educativa.
No entanto, pontua -se aqui, algumas \textbf{diretrizes} legais e/ou de caráter orientador.
O \textbf{Parecer} do CNE/CEB No 20/2009, orienta que o número de crianças por professor deve possibilitar atenção, responsabilidade e interação com as crianças e suas famílias e que, ao considerar a dimensão do espaço físico e das crianças, em se tratando de agrupamentos com crianças de mesma faixa de idade, a recomendação é de que obedeça a proporção de 6 a 8 crianças por professor ( no caso de crianças até um ano ), 15 crianças por professor ( no caso de crianças de dois e três anos ) e 20 crianças por professor ( nos agrupamentos de crianças de quatro e cinco anos ).
Neste mesmo sentido, a Resolução CEE no 01/2016 indica que a \textbf{instituição} de ensino adotará a organização de agrupamento de seis a oito crianças por professor, no caso de crianças de até um ano de idade ; de 15 crianças por professor, no caso de crianças de dois e três anos de idade ; de 20 a 25 crianças por professor, no caso de crianças de quatro e cinco anos.
A referida Resolução destaca ainda, que o número de crianças em cada agrupamento deverá possibilitar locomoção e interação das crianças entre elas, suas famílias e professores, considerando as características do espaço físico e o desenvolvimento infantil.
Além disso, afirma que não deve haver agrupamentos de crianças da Educação Infantil com crianças do Ensino Fundamental em uma mesma turma.
A respeito da \textbf{gestão} democrática, os Parâmetros Nacionais de Qualidade para a Educação Infantil ( BRASIL, 2006 ), em consonância com os princípios da LDB, abordam que, para definição e implementação de \textbf{políticas} para essa etapa educativa, cabe às \textbf{Secretarias} \textbf{Municipais} de Educação, a \textbf{garantia} da \textbf{gestão} democrática com a implantação de conselhos nas \textbf{instituições} públicas, de forma a haver o aprimoramento das formas de participação da comunidade.
Sobre este aspecto, ainda se constitui como um desafio para as \textbf{instituições} de Educação Infantil em nosso estado, o alcance dessa experiência de escolher seu representante por meio de eleições diretas para gestores educacionais.
No entanto, são reconhecidos os esforços nesse sentido, e é sabido que, é possível assumir alguns princípios de \textbf{gestão} democrática e participativa, mesmo sem esse mecanismo legitimo de seleção de gestores.
Busca -se ao longo da história estabelecer a ideia de descentralizar a \textbf{gestão} promovendo maior participação da sociedade, por meio de práticas democráticas.
Pois, se evidencia a importância da participação da comunidade nas \textbf{instituições} de educação.
O contexto educacional contemporâneo exige uma \textbf{gestão} escolar participativa e democrática e a função do gestor, ao longo da história, vem se modificando, na medida em que a sociedade muda e se transforma, demandando novas posturas e mecanismos que superam a \textbf{gestão} centralizadora e burocrática.
De acordo com Veiga ( 1997 ), tal pensamento rompe com a separação entre concepção e execução, entre o pensar e o fazer, entre a teoria e a prática.
Nesse sentido, a \textbf{instituição} democrática fomenta o exercício da democracia a partir de uma \textbf{gestão} compartilhada, participativa e compromissada com a efetivação da descentralização de decisões, impulsionando crianças, professores, \textbf{funcionários}, famílias, Conselhos Escolar e Fiscal e toda comunidade escolar a vivências que propiciem a participação ativa dos pares.
Ressalta -se nesse documento, o necessário movimento democrático nas \textbf{instituições} de Educação Infantil de garantir a constituição de seus conselhos escolares e uma \textbf{gestão} articulada com a comunidade educativa.
Reforça -se, sobretudo, uma boa relação de parceria com as famílias e outros setores comunitários e sociais.
O \textbf{Parecer} CNE/CEB No 20/2009 afirma que a Educação Infantil precisa refletir sobre o modo como deve integrar as ações e os projetos educacionais às famílias, devendo ser mantida e desenvolvida essa integração durante todo o tempo em que as crianças estejam na \textbf{Instituição}.
Aborda, portanto, aspectos relacionados à participação da família na \textbf{gestão} da proposta pedagógica, junto aos professores e demais profissionais da educação, nos conselhos escolares, no acompanhamento de projetos e nas atividades promovidas pela \textbf{instituição}, a fim de contribuir com o desenvolvimento das crianças.
Para que a participação se concretize, é necessário que os pais sejam ouvidos enquanto representantes das crianças, principalmente as que ainda são muito pequenas.
As DCNEI ( BRASIL, 2009a ) reafirmam que as \textbf{instituições} devem assumir a responsabilidade de compartilhar e complementar a educação e o cuidado das crianças com as famílias e apontam que deve se assegurar a participação, o diálogo e a escuta cotidiana das famílias, bem como o respeito e a valorização de suas formas de organização.
Nessa mesma perspectiva, a BNCC ( BRASIL, 2017 ) destaca que o diálogo entre a \textbf{instituição} de Educação Infantil e a família, bem como o compartilhamento de responsabilidades, são modos essenciais que potencializam as aprendizagens e o desenvolvimento das mesmas.
Levando em consideração tais orientações, é importante que a \textbf{instituição} tenha clareza da relevância do seu papel como mobilizadora de ações que impulsionem o envolvimento da família, considerando a realidade de sua comunidade local, suas limitações e disponibilidades.
Nesse sentido, uma atenção especial deve ser dada ao conhecimento das particularidades de cada família, por exemplo, quanto às suas atividades profissionais, que podem ser as razões causadoras da não participação nas atividades promovidas pela \textbf{instituição}, dentre outros fatores sociais.
Assim, a \textbf{instituição} pode redefinir suas propostas de acolhimento às famílias, adequando -as às suas possibilidades e desse modo, conseguindo inseri -las no processo participativo a que se propõe.
Com isso, permite seu envolvimento no projeto curricular, de modo que estas contribuam com o trabalho desenvolvido, sentindo -se corresponsável e reconhecendo na \textbf{instituição} educativa, um lugar de transformação social.
Nas \textbf{instituições} educativas de nosso Estado, são desenvolvidas estratégias importantes de articulação com as famílias, como as reuniões pedagógicas mensais, visitas domiciliares, participações em eventos, reuniões bimestrais, palestras, seminários, dentre outros.
Ressalta -se aqui, o necessário \textbf{fortalecimento} dessas ações na constituição de uma \textbf{política} efetiva, importante e essencial, nos desdobramentos dos processos pedagógicos e ou administrativos na \textbf{instituição} educativa envolvendo a família e a comunidade.
Além da necessária integração da \textbf{Instituição} com as famílias, outras mediações podem ser viabilizadas com as \textbf{instituições} de controle social responsáveis pelo acompanhamento das \textbf{políticas} sociais em cada \textbf{município}.
Nesse contexto, os conselhos escolares têm um papel fundamental na \textbf{garantia} desse processo de participação social.
Articulado aos princípios da \textbf{gestão} e ao processo pedagógico fomentado na \textbf{Instituição} de Educação infantil encontra -se o Coordenador Pedagógico que é o mediador das demandas pedagógicas e formativas.
Em nosso estado, a função de coordenação pedagógica é assumida por professores que, em alguns casos, são responsáveis por uma \textbf{instituição} ou pelo atendimento a todos os professores da Rede.
Assim como \textbf{compete} à \textbf{instituição} a aprendizagem e o desenvolvimento das crianças, é preciso também garantir espaços e tempos formativos, articulados pelo coordenador pedagógico, que oportunizem a aprendizagem dos professores, conforme destaca Libâneo ( 2009, p.11 ) : > O aprendizado cotidiano nas atividades conjuntas de colaboração entre professores precisa ser reforçado por outras atividades de formação continuada no contexto de trabalho.
São aquelas práticas formativas de desenvolvimento pessoal e profissional do pessoal da escola : grupos de estudo, projetos de trabalho, encontros pedagógicos para troca de ideias e experiências, as entrevistas com a Coordenação Pedagógica, etc.
Sobre esse assunto, a Resolução no 01/2016 – CEE/RN, de 13 de abril de 2016, destaca o aperfeiçoamento e a atualização dos professores, que deve ocorrer através das \textbf{instituições} mantenedoras da Educação Infantil.
Visa, portanto, assegurar que a formação continuada esteja em coerência com a proposta pedagógica e as \textbf{políticas} relacionadas a essa etapa educacional.
Na mesma perspectiva, o \textbf{Parecer} CNE/CEB no : 20/2009 compreende os programas de formação como direito dos professores, visto que contribuem para a melhoria da prática e para o seu desenvolvimento como pessoa e como profissional no exercício da atividade docente.
Também são enfatizados na lei os requisitos básicos para formação que contribuam com uma Educação Infantil de qualidade.
Sendo assim, considera -se os aspectos abordados neste \textbf{capítulo}, como proposições a serem pensadas e materializadas no desenvolvimento dos projetos políticos pedagógicos em cada \textbf{instituição} educativa.

% matched lemmas: artigo, capítulo, competir, diretriz, fortalecimento, funcionário, garantia, gestão, instituição, municipal, município, normatização, parecer, permanecer, política, regulamentação, secretaria
\end{textsample}
